% !TeX root = ../../Book1.tex
% 纲要标题：## 第7章　正规列、可解群与幂零群
% 纲要位置：计划纲要.md，第 300 行。

\chapter{正规列、可解群与幂零群}
\label{b1:ch:07}
\BookChapterNumber{b1:ch:07}

\begin{chapterabstract}
群可以沿正规子群逐层分解，也可以用交换子逐步消去非交换性。本章证明合成因子的唯一性，建立可解群和幂零群的基本判据，并证明五次以上交错群的单性。这些结果把前面零散的有限群算例组织成结构理论，同时为后面的 Galois 理论准备可解性这一判别工具。
\end{chapterabstract}

\begin{learninggoals}
\begin{itemize}
\item 分清正规列与次正规列，掌握蝴蝶引理如何配对细化的因子。
\item 会计算常见群的导出列和中心列，并检查可解性、幂零性的稳定性质。
\item 能把有限幂零群拆为 Sylow 子群直积，理解交错群单性证明中的最小支持度方法。
\end{itemize}
\end{learninggoals}

本章使用\chapref{b1:ch:02}的同构定理、\chapref{b1:ch:04}的类方程及\chapref{b1:ch:05}的 Sylow 理论。\secref{b1:sec:0701}的合成列理论与后面的交换子方法可先分别阅读，再通过“有限可解群的合成因子”联系起来。伴读可选 Dummit 与 Foote 的《Abstract Algebra》\cite{DummitFoote2003}中群结构的有关内容；本章所用的结构定理均在正文证明。

% !TeX root = ../../Book1.tex
% 纲要标题：### 7.1 正规列与合成列
% 纲要位置：计划纲要.md，第 302 行。

\section{正规列与合成列}
\label{b1:sec:0701}

% 纲要内容（仅作写作计划，尚非教材正文）：
%
% * 次正规列、合成列及细化
% * Schreier 细化定理
% * Jordan–Hölder 定理
%

\subsection{次正规列、合成列及细化}
\label{b1:subsec:0701:01}

前两章用一个正规子群把群分成子群与商群。反复进行这种分解，便得到一列较简单的因子。不过，某个群可以有许多不同的正规子群；要使分解成为结构工具，必须说明不同选择之间有什么共同之处。

\begin{definition}\label{b1:def:subnormal-series}
群 \(G\) 的一个\term[cizhenggui-lie]{次正规列}[subnormal series]是有限链
\[
G=G_0\trianglerighteq G_1\trianglerighteq\cdots
\trianglerighteq G_r=1,
\]
其中只要求 \(G_{i+1}\) 在 \(G_i\) 中正规。若各 \(G_i\) 都在 \(G\) 中正规，称为\term[zhenggui-lie]{正规列}[normal series]。商群 \(G_i/G_{i+1}\) 称为列的因子；插入若干子群得到的新次正规列称为原列的\term[xihua]{细化}[refinement]。
\end{definition}
有些文献把本书的次正规列也称为正规列。这里区分二者，是因为“在前一项中正规”不推出“在整个群中正规”。例如
\[
S_4\trianglerighteq V_4\trianglerighteq
\langle(12)(34)\rangle\trianglerighteq1
\]
是次正规列，但中间的二阶子群不在 \(S_4\) 中正规。

\begin{definition}\label{b1:def:composition-series}
回顾单群是只有自身和单位子群两个正规子群的非平凡群。所有因子都是单群的严格次正规列称为\term[hecheng-lie]{合成列}[composition series]，其因子称为\term[hecheng-yinzi]{合成因子}[composition factor]。
\end{definition}
对应定理说明，合成列恰是不能再插入不同子群的严格次正规列。每个非平凡有限群都存在合成列：选一个极大真正规子群，再对它重复；子群的阶严格下降，过程必停。无限群未必如此，\(\mathbb Z\) 的任何非零子群都同构于 \(\mathbb Z\)，都还含有非零真子群，因而不可能在有限步后到达一个单的末项。

\ProseParagraphBreak
若删去平凡因子后，可以把两列的因子一一配对，使每对同构，就称它们为\term[dengjia-lie]{等价列}[equivalent series]；配对不要求保留原有次序。例如 \(C_6\) 可以先取二阶子群，也可以先取三阶子群，两次所得因子分别是 \(C_3,C_2\) 与 \(C_2,C_3\)。

\subsection{Schreier 细化定理}
\label{b1:subsec:0701:02}

比较两条列时，应把一条列中的子群逐个与另一条列相交，再乘回原来较低的一项。使这些商群成对同构的工具是\term[hudie-yinli]{蝴蝶引理}[butterfly lemma]，也称 Zassenhaus 引理。

\begin{lemma}[Zassenhaus]\label{b1:lem:butterfly}
设 \(A'\trianglelefteq A\)、\(B'\trianglelefteq B\) 是同一群中的子群。则
\[
\frac{A'(A\cap B)}{A'(A\cap B')}
\cong
\frac{B'(B\cap A)}{B'(B\cap A')}.
\]
式中分母分别正规于分子。
\end{lemma}
\begin{proof}
令 \(C=A\cap B\)，以及
\[
D=(A'\cap B)(A\cap B').
\]
两个乘积因子都正规于 \(C\)，所以 \(D\trianglelefteq C\)。记
\(K=A'(A\cap B)\)、\(L=A'(A\cap B')\)。映射 \(K\to A/A'\) 的像由 \(C\) 给出；\(L\) 的像是 \(A\cap B'\) 的像。由于 \(A\cap B'\) 正规于 \(C\)，这个像正规于 \(K/A'\)，从而 \(L\trianglelefteq K\)。

包含映射 \(C\to K\) 后接商映射给出满同态 \(C\to K/L\)。其核为 \(C\cap L=D\)：若 \(c=a'b'\in C\)，其中 \(a'\in A'\)、\(b'\in A\cap B'\)，则 \(a'=c(b')^{-1}\in A'\cap B\)；反向包含直接成立。于是 \(K/L\cong C/D\)。交换 \((A,A')\) 与 \((B,B')\)，另一个商群也同构于 \(C/D\)。
\end{proof}

\begin{theorem}[Schreier]\label{b1:thm:schreier-refinement}
同一群的任意两条次正规列都有等价的细化。
\end{theorem}
\begin{proof}
把两条列记为 \((G_i)_{i=0}^r\) 与 \((H_j)_{j=0}^s\)。在 \(G_i\) 和 \(G_{i+1}\) 之间插入
\[
G_{ij}=G_{i+1}(G_i\cap H_j),\qquad 0\le j\le s.
\]
它的首末项分别为 \(G_i\)、\(G_{i+1}\)。蝴蝶引理保证相邻项的正规性，因此把这些小段接起来得到第一条列的细化。同样在 \(H_j\) 与 \(H_{j+1}\) 之间插入
\[
H_{ji}=H_{j+1}(H_j\cap G_i),\qquad 0\le i\le r.
\]
对每对 \((i,j)\)，蝴蝶引理给出
\[
G_{ij}/G_{i,j+1}\cong H_{ji}/H_{j,i+1}.
\]
两个矩形数组按行和按列读取的次序不同，但格子的对应提供了所需配对。重复项只产生平凡因子，删掉即可。
\end{proof}
这就是\term[schreier-xihua-dingli]{Schreier 细化定理}[Schreier refinement theorem]。它比较的是商群，而不是断言两条细化中的子群相同。事实上，\(C_2\times C_2\) 中任取一个二阶子群都给出合成列，三种选择的中间子群互不相同。

\subsection{Jordan–Hölder 定理}
\label{b1:subsec:0701:03}

\begin{theorem}[Jordan--Hölder]\label{b1:thm:jordan-holder}
若群 \(G\) 有合成列，则它的任意两条合成列长度相同，并且合成因子在同构及排列的意义下相同。
\end{theorem}
\begin{proof}
对两条合成列应用 Schreier 细化定理。合成列中的相邻因子是单群；由对应定理，在相邻两项之间不能插入不同的新子群。因此细化只会重复原有项。删去这些重复项之后，细化的非平凡因子就是原列的全部因子。两条细化等价，结论随即成立。
\end{proof}
\term[jordan-holder-dingli]{Jordan--Hölder 定理}[Jordan--Hölder theorem]使合成因子的无序集合（计入重数）成为群的不变量。不过这些因子不足以重建原群：\(C_4\) 和 \(C_2\times C_2\) 都有两个 \(C_2\) 因子，却不相互同构。合成因子记录了分解后的简单部分，部分之间怎样结合仍是扩张问题。

\ProseParagraphBreak
以 \(S_4\) 为例，链
\[
S_4\trianglerighteq A_4\trianglerighteq V_4
\trianglerighteq\langle(12)(34)\rangle\trianglerighteq1
\]
给出因子 \(C_2,C_3,C_2,C_2\)，故是合成列。因子阶的乘积为 \(24\)，与 \(S_4\) 的阶相符。以后遇到一个有限群的候选合成列时，可以先用阶数乘积排查计算错误，再逐项检查正规性和单性；阶数一致本身并不能代替后两项检查。

% !TeX root = ../../Book1.tex
% 纲要标题：### 7.2 可解群
% 纲要位置：计划纲要.md，第 308 行。

\section{可解群}
\label{b1:sec:0702}

% TODO: 按以下纲要要点撰写本节。

% 纲要内容（仅作写作计划，尚非教材正文）：
%
% * 导出列与交换子
% * 可解性的子群、商群及扩张稳定性
% * 对称群和三角矩阵群的例子
%

\subsection{导出列与交换子}
\label{b1:subsec:0702:01}

待撰写。

\subsection{可解性的稳定性}
\label{b1:subsec:0702:02}

待撰写。

\subsection{对称群与三角矩阵群}
\label{b1:subsec:0702:03}

待撰写。

% !TeX root = ../../Book1.tex
% 纲要标题：### 7.3 幂零群
% 纲要位置：计划纲要.md，第 314 行。

\section{幂零群}
\label{b1:sec:0703}

% 纲要内容（仅作写作计划，尚非教材正文）：
%
% * 上、下中心列
% * 有限 p 群的幂零性
% * 有限幂零群与 Sylow 子群直积
%

\subsection{上、下中心列}
\label{b1:subsec:0703:01}

可解群允许每一步的商群交换，却不要求该商群与群中其他元素相容。若一层元素不仅彼此交换，而且在模去更低一层以后与整个群交换，就得到更强的条件。

\begin{definition}\label{b1:def:central-series}
群 \(G\) 的\term[xia-zhongxin-lie]{下中心列}[lower central series]定义为
\[
\gamma_1(G)=G,\qquad \gamma_{i+1}(G)=[G,\gamma_i(G)].
\]
其\term[shang-zhongxin-lie]{上中心列}[upper central series]定义为
\[
Z_0(G)=1,\qquad
Z_{i+1}(G)/Z_i(G)=Z\bigl(G/Z_i(G)\bigr).
\]
右式通过商映射的原像定义 \(Z_{i+1}(G)\)。若 \(\gamma_{c+1}(G)=1\)，称 \(G\) 为\term[mi-ling-qun]{幂零群}[nilpotent group]；最小这样的非负整数 \(c\) 是其幂零类。
\end{definition}
\symidx{\(\gamma_i(G)\)：下中心列}
\symidx{\(Z_i(G)\)：上中心列}
这些子群都是特征子群。对下中心列，这由交换子在自同构下的相容性归纳得出；对上中心列，自同构先诱导商群的自同构，再保持该商群的中心。上中心列的定义也可写为
\[
x\in Z_{i+1}(G)\quad\Longleftrightarrow\quad
[g,x]\in Z_i(G)\text{ 对每个 }g\in G\text{ 成立}.
\]

\begin{proposition}\label{b1:prop:central-series-equivalence}
对任意 \(c\ge0\)，以下条件等价：
\begin{enumerate}
\item \(\gamma_{c+1}(G)=1\)；
\item \(Z_c(G)=G\)；
\item 存在正规子群链 \(1=C_0\le C_1\le\cdots\le C_c=G\)，使 \([G,C_i]\le C_{i-1}\)。
\end{enumerate}
每个幂零群都可解。
\end{proposition}
\begin{proof}
第一条件成立时，取 \(C_i=\gamma_{c+1-i}(G)\)，即得第三条件。若第三条件成立，由 \(C_0=Z_0(G)\) 和上述中心的判别逐步得到 \(C_i\le Z_i(G)\)，所以第二条件成立。第二条件成立时，直接取 \(C_i=Z_i(G)\) 得到第三条件；由 \(\gamma_1(G)=C_c\) 及交换子的单调性可归纳推出 \(\gamma_j(G)\le C_{c+1-j}\)，于是得到第一条件。

第三条件的每个因子 \(C_i/C_{i-1}\) 都阿贝尔，因此 \(G\) 可解。
\end{proof}
非平凡阿贝尔群的幂零类为一。四元数群满足
\[
Q_8'=Z(Q_8)=\{1,-1\},\qquad [Q_8,\{1,-1\}]=1,
\]
所以幂零类为二。相反，\(S_3\) 虽然可解，中心却平凡；其上中心列永远停留在单位群，故不幂零。由交换子的包含关系还可得幂零性传给子群及商群，但“可解扩张仍可解”不能换成“幂零扩张仍幂零”：\(C_3\trianglelefteq S_3\) 的核与商都阿贝尔就是反例。

\subsection{有限 \texorpdfstring{\(p\)}{p} 群的幂零性}
\label{b1:subsec:0703:02}

\begin{theorem}\label{b1:thm:p-group-nilpotent}
每个有限 \(p\) 群都幂零。若群的阶为 \(p^a\)，则幂零类至多为 \(a\)。
\end{theorem}
\begin{proof}
\chapref{b1:ch:04}的类方程已经说明，非平凡有限 \(p\) 群的中心非平凡。如果 \(Z_i(G)\ne G\)，那么 \(G/Z_i(G)\) 仍是非平凡有限 \(p\) 群，故其中心非平凡，得到严格包含 \(Z_i(G)<Z_{i+1}(G)\)。每次严格增大至少使阶乘上 \(p\)；从一阶增大到 \(p^a\) 至多需要 \(a\) 步。因此 \(Z_a(G)=G\)。
\end{proof}
这里并非只使用“中心非平凡”一次，而是在每一层商群中重新应用它。一般非平凡中心并不足以保证幂零，例如 \(C_2\times S_3\) 的中心非平凡，但其商群含有不可幂零的 \(S_3\)。

\ProseParagraphBreak
对阶为 \(8\) 的二面体群 \(D_8=\langle r,s\mid r^4=s^2=1,\ srs^{-1}=r^{-1}\rangle\)，有
\[
[s,r]=r^{-2}=r^2,\quad
\gamma_2(D_8)=\langle r^2\rangle,\quad
\gamma_3(D_8)=1.
\]
它与 \(Q_8\) 同为二类幂零群，但两群不相互同构：前者有五个二阶元，后者只有一个。幂零类衡量的是交换子消失的层数，不是群的完全分类资料。

\subsection{有限幂零群与 Sylow 子群直积}
\label{b1:subsec:0703:03}

有限群的幂零性还可以从 Sylow 子群是否彼此独立看出来。证明先用一个对无限幂零群也成立的观察。

\begin{lemma}\label{b1:lem:nilpotent-normalizer}
幂零群的每个真子群 \(H<G\) 都严格包含在其正规化子 \(N_G(H)\) 中。
\end{lemma}
\begin{proof}
取最小的 \(i\) 使 \(Z_i(G)\not\le H\)，再取 \(z\in Z_i(G)\setminus H\)。此时 \(Z_{i-1}(G)\le H\)。对任意 \(h\in H\)，有 \([z,h]\in Z_{i-1}(G)\le H\)，于是 \(zhz^{-1}\in H\)。对 \(z^{-1}\in Z_i(G)\) 作同样论证，得到 \(zHz^{-1}=H\)。所以 \(z\in N_G(H)\setminus H\)。
\end{proof}

\begin{theorem}\label{b1:thm:finite-nilpotent-sylow}
对有限群 \(G\)，以下条件等价：\(G\) 幂零；\(G\) 的每个 Sylow 子群都正规；\(G\) 是其各个 Sylow 子群的内直积。
\end{theorem}
\begin{proof}
设 \(G\) 幂零，\(P\) 是一个 Sylow 子群，令 \(N=N_G(P)\)。在 \(N\) 中，\(P\) 是正规 Sylow 子群，故是唯一的对应素数的 Sylow 子群，从而是 \(N\) 的特征子群。任何正规化 \(N\) 的元素都保持 \(P\)，所以
\(N_G(N)\le N_G(P)=N\)。若 \(N<G\)，这与上一引理矛盾。因此 \(N=G\)，\(P\trianglelefteq G\)。

反过来，设各 Sylow 子群 \(P_1,\ldots,P_t\) 正规。对 \(i\ne j\)，正规性给出 \([P_i,P_j]\le P_i\cap P_j=1\)，所以不同因子逐元素交换。乘积映射
\[
P_1\times\cdots\times P_t\longrightarrow G,
\qquad (x_1,\ldots,x_t)\longmapsto x_1\cdots x_t
\]
是同态。其核平凡：若某一 \(x_i\) 等于其余因子乘积的逆，那么其阶一方面是对应素数的幂，另一方面整除其余因子群阶的乘积，只能为一。定义域阶等于 \(|G|\)，故此映射为同构。

最后，有限 \(p\) 群已经证明幂零，而直积的交换子逐分量计算，因此
\(\gamma_k(P_1\times\cdots\times P_t)=\prod_i\gamma_k(P_i)\)。取各因子幂零类的最大值，得到直积幂零。
\end{proof}
例如，阶为 \(12\) 的 \(C_3\times C_4\) 和 \(C_3\times V_4\) 均幂零，\(A_4\) 却不幂零，因为其 Sylow 三子群不唯一。判断有限群时，可先检查 Sylow 子群；计算无限群时，则应回到中心列的定义，不能直接搬用有限 Sylow 直积判据。

% !TeX root = ../../Book1.tex
% 纲要标题：### 7.4 交错群的单性
% 纲要位置：计划纲要.md，第 320 行。

\section{交错群的单性}
\label{b1:sec:0704}

% 纲要内容（仅作写作计划，尚非教材正文）：
%
% * 三轮换生成与正规子群分析
% * n≥5 时交错群的单性
% * 与一般方程不可根式求解的后续联系
%
% ---
%

\subsection{三轮换生成与正规子群}
\label{b1:subsec:0704:01}

一个非平凡正规子群一旦包含某种置换，就必须包含它的所有群内共轭。因此，证明交错群单性的关键是找到足够简单且共轭丰富的元素。三轮换最适合这个任务：它是移动字母数最少的非单位偶置换。

\begin{lemma}\label{b1:lem:three-cycles-generate}
三轮换生成 \(A_n\)。当 \(n\ge5\) 时，所有三轮换在 \(A_n\) 中共轭。
\end{lemma}
\begin{proof}
三轮换生成 \(A_n\) 已在\secref{b1:sec:0302}证明；这里只需补充群内共轭性。

任意两个三轮换 \(\alpha,\beta\) 在 \(S_n\) 中共轭，设 \(u\alpha u^{-1}=\beta\)。若 \(u\) 为偶置换，已经得到所需共轭。若 \(u\) 为奇置换，因 \(n\ge5\)，可在 \(\beta\) 的支集之外选两个字母构成换位 \(t\)。它与 \(\beta\) 交换，且 \(tu\) 为偶置换，于是 \((tu)\alpha(tu)^{-1}=\beta\)。
\end{proof}
因此只要证明非平凡正规子群包含一个三轮换，就能推出它等于整个 \(A_n\)。下面的交换子方法把任意非单位元改造成移动较少字母的元素。

\begin{definition}\label{b1:def:permutation-support}
置换 \(\sigma\) 的\term[zhichi]{支集}[support]是
\(\operatorname{supp}(\sigma)=\{i\mid\sigma(i)\ne i\}\)。其大小称为此处的支持度。
\end{definition}
若 \(N\trianglelefteq A_n\)、\(\sigma\in N\)、\(\tau\in A_n\)，则
\[
[\sigma,\tau]=\sigma\tau\sigma^{-1}\tau^{-1}\in N.
\]
如果 \(\tau\) 是三轮换，两个乘积因子的支集各有三个字母，故交换子的支持度至多为六。还必须确保交换子非平凡；下述证明在每次使用时都核对这一点。

\subsection{交错群的单性}
\label{b1:subsec:0704:02}

\begin{theorem}\label{b1:thm:alternating-simple}
当 \(n\ge5\) 时，交错群 \(A_n\) 是非阿贝尔单群。
\end{theorem}
\begin{proof}
设 \(1\ne N\trianglelefteq A_n\)，在 \(N\setminus\{1\}\) 中取支持度最小的元素 \(\sigma\)，记该支持度为 \(m\)。偶置换不可能恰好移动一个或两个字母，所以 \(m\ge3\)。

先证 \(m\le6\)。若 \(m>6\)，选 \(a\) 满足 \(\sigma(a)=b\ne a\)，再选不同于 \(a,b\) 的两个字母 \(c,d\)，令 \(\tau=(acd)\)。共轭 \(\sigma\tau\sigma^{-1}\) 的支集包含 \(b\)，而 \(\tau\) 的支集不包含 \(b\)，故二者不同，\([\sigma,\tau]\ne1\)。这个交换子属于 \(N\)，支持度至多六，违背 \(m\) 的最小性。

若 \(m=3\)，则 \(\sigma\) 已是三轮换。其余可能的偶置换轮换类型只有下面几种；字母均表示互不相同的点。
\begin{enumerate}
\item 若 \(m=4\)，必有 \(\sigma=(ab)(cd)\)。选支集外的字母 \(e\)（此处用 \(n\ge5\)），令 \(\tau=(cde)\)。于是
\(\sigma\tau\sigma^{-1}=(dce)=\tau^{-1}\)，所以
\([\sigma,\tau]=\tau^{-2}=\tau\) 是 \(N\) 中的三轮换。这也表明支持度最小元不能实际具有 \(m=4\)。
\item 若 \(m=5\)，偶性迫使 \(\sigma=(abcde)\)。取 \(\tau=(abc)\)，则共轭为 \((bcd)\)。二者支集不同，所以交换子非平凡；其支集包含于 \(\{a,b,c,d\}\)，违背最小性。
\item 若 \(m=6\) 且 \(\sigma=(abc)(def)\)，取 \(\tau=(abd)\)。共轭为 \((bce)\)，所以交换子非平凡且支集包含于 \(\{a,b,c,d,e\}\)，仍与最小性矛盾。
\item 若 \(m=6\) 且 \(\sigma=(abcd)(ef)\)，取 \(\tau=(abc)\)。共轭为 \((bcd)\)，产生一个非平凡且支持度至多四的交换子，再次矛盾。
\end{enumerate}
这些情形穷尽了支持度四至六的偶置换：六轮换、三个换位，以及三轮换乘一个换位都是奇置换；含固定点的轮换型应按实际支持度计入前面的情形。因此 \(N\) 含有三轮换。由前一引理及正规性，\(N\) 包含所有三轮换，故 \(N=A_n\)。最后，\((123)\) 与 \((124)\) 不交换，说明 \(A_n\) 非阿贝尔。
\end{proof}
\term[jiaocuo-qun-danxing]{交错群的单性}[simplicity of the alternating group]在 \(n=5\) 就成立；证明没有使用多于五个字母去处理 \(A_5\) 的双换位。低阶例外也应同时记住：\(A_3\cong C_3\) 是阿贝尔单群，\(A_4\) 有真正规子群 \(V_4\)，因而不单。

\subsection{交错群与根式不可解性}
\label{b1:subsec:0704:03}

\begin{corollary}\label{b1:cor:symmetric-nonsolvable}
当 \(n\ge5\) 时，\(A_n'=A_n\)，因而 \(A_n\) 及 \(S_n\) 均不可解。对 \(n\le4\)，\(S_n\) 可解。
\end{corollary}
\begin{proof}
导出子群正规。由于 \(A_n\) 非阿贝尔，其导出子群不平凡；单性迫使它等于 \(A_n\)。所以导出列不可能下降到单位群，\(A_n\) 不可解，而可解性传给子群表明 \(S_n\) 也不可解。低阶情形在\secref{b1:sec:0702}已经算出；\(S_1,S_2\) 则直接阿贝尔。
\end{proof}
这为一般方程的根式不可解性准备了群论部分。到本卷 Galois 理论中，我们将建立“用有限次开根表达根”与适当 Galois 群的可解性之间的联系，并证明相应一般多项式的 Galois 群是对称群。两项桥梁建立之后，\(S_n\) 在 \(n\ge5\) 时不可解才转化为一般五次及更高次方程没有统一根式公式。

\ProseParagraphBreak
当前结果并不说“每个五次方程都不可根式求解”。例如 \(x^5-2=0\) 的一个根本来就是 \(\sqrt[5]{2}\)，其余根再乘五次单位根即可表达。决定具体方程的是它自己的 Galois 群，而不是次数一个数字。这里先明确逻辑分工：本章完成可解群及非阿贝尔单群的事实，域扩张与根式之间的对应留到本卷后续章节证明。


\begin{chaptersummary}
Schreier 细化定理通过交与积比较不同的次正规列，Jordan--Hölder 定理则把合成因子确定为计入重数的群不变量。导出列检测群能否由有限层交换群逐步扩张得到；中心列提出更强要求，使每一层在相应商群中居于中心。有限幂零群恰是各 Sylow 子群的直积，但可解群一般不具备这种分离性质。\(A_n\) 在 \(n\ge5\) 时的非交换性与单性给出了导出列永不下降的基本例子。
\end{chaptersummary}
\BookExercises

\begin{exercise}\label{b1:ex:07-composition-dihedral}
设 \(D_{24}=\langle r,s\rangle\)，其中 \(r\) 的阶为 \(12\)。写出一条合成列，确定合成因子，并说明 \(C_{24}\) 与它有相同合成因子但不相互同构。
\end{exercise}
\begin{solution}
取
\[
D_{24}\trianglerighteq\langle r\rangle
\trianglerighteq\langle r^2\rangle
\trianglerighteq\langle r^4\rangle\trianglerighteq1.
\]
\cref{b1:prop:power-order}给出 \(r^2,r^4\) 的阶分别为 \(6,3\)，所以子群阶依次为 \(24,12,6,3,1\)。旋转子群的指数为二，由\cref{b1:prop:normal-basic-examples}正规；后续相邻项处于循环群内，而循环群是交换群，故也满足正规性。\cref{b1:thm:lagrange}给出相邻商的阶为 \(2,2,2,3\)，由\cref{b1:prop:abelian-simple}及素数阶群的循环性，它们分别是单群 \(C_2,C_2,C_2,C_3\)，因此所写链是合成列。\cref{b1:thm:cyclic-subgroups}保证 \(C_{24}\) 中存在阶为 \(12,6,3,1\) 的嵌套子群，取此列得到相同因子。前者非交换，因为 \(srs^{-1}=r^{-1}\ne r\)；后者交换，所以不相互同构。
\end{solution}

\begin{exercise}\label{b1:ex:07-derived-map}
设 \(f:G\to H\) 是满同态，\(N=\ker f\)。证明 \(H'\cong G'/(G'\cap N)\)。若 \(G'\cap N=1\) 且 \(H\) 交换，能否推出 \(G\) 交换？
\end{exercise}
\begin{solution}
由式\eqref{b1:eq:commutator-image}及 \(f(G)=H\)，限制同态 \(f|_{G'}:G'\to H'\) 的像是整个 \(H'\)，其核为 \(G'\cap\ker f=G'\cap N\)。对这个满同态应用\cref{b1:thm:first-isomorphism}，即得所求同构。若此外 \(H'=1\)、\(G'\cap N=1\)，则 \(G'\) 既包含在 \(N\) 中又与 \(N\) 交为一，因此 \(G'=1\)。\cref{b1:prop:derived-quotient}取正规子群为 \(1\)，便说明 \(G\) 交换。
\end{solution}

\begin{exercise}\label{b1:ex:07-normal-solvable-product}
设 \(M,N\trianglelefteq G\) 都可解。证明 \(MN\) 可解，并据此说明每个有限群有唯一的最大可解正规子群。
\end{exercise}
\begin{solution}
\cref{b1:thm:second-isomorphism}用于 \(M\trianglelefteq G\)、\(N\le G\)，给出 \(MN\) 是子群，\(M\trianglelefteq MN\)，以及 \(MN/M\cong N/(M\cap N)\)；这个同构把 \(n(M\cap N)\) 送到 \(nM\)。由\cref{b1:thm:solvable-closure}的商群结论，右侧作为 \(N\) 的商群可解；再将同一定理的扩张结论用于 \(M\trianglelefteq MN\)，得到 \(MN\) 可解。又因 \(M,N\) 均在 \(G\) 中正规，任意共轭保持它们的乘积，故 \(MN\trianglelefteq G\)。有限群只有有限多个子群，把它的全部可解正规子群依次相乘，反复应用刚证的两子群乘积结论，得到可解正规子群 \(R\)。它包含每个可解正规子群，所以最大；任何另一个具有此性质的子群必须既包含 \(R\) 又被 \(R\) 包含，故唯一。
\end{solution}

\begin{exercise}\label{b1:ex:07-dihedral-series}
对 \(n\ge3\)，计算 \(D_{2n}'\) 及下中心列，并证明 \(D_{2n}\) 幂零当且仅当 \(n\) 为二的幂。若 \(n=2^a\)，求其幂零类。
\end{exercise}
\begin{solution}
由\cref{b1:prop:dihedral-normal-form}中的 \(srs^{-1}=r^{-1}\)，计算得 \([s,r]=r^{-2}\)，故 \(\langle r^2\rangle\le D_{2n}'\)。共轭生成元 \(r,s\) 都保持 \(\langle r^2\rangle\)，所以这个子群正规。模去它以后，\(r\) 的像满足 \(r=r^{-1}\)，因而与 \(s\) 的像交换，商群交换；\cref{b1:prop:derived-quotient}给出反向包含，得到 \(D_{2n}'=\langle r^2\rangle\)。对任何整数 \(t,k\)，标准形 \(r^t\) 与 \(r^ts\) 分别满足
\[
[r^t,r^k]=1,\qquad [r^ts,r^k]=r^{-2k}.
\]
因此与 \(\langle r^k\rangle\) 取交换子恰生成 \(\langle r^{2k}\rangle\)，对下中心列的指标归纳即得
\[
\gamma_j(D_{2n})=\langle r^{2^{j-1}}\rangle\qquad(j\ge2).
\]
\cref{b1:prop:power-order}说明 \(r^{2^{j-1}}=1\) 当且仅当 \(n\mid2^{j-1}\)，所以该列终止当且仅当 \(n\) 是二的幂。若 \(n=2^a\)（\(a\ge2\)），则 \(\gamma_{a+1}=1\)、\(\gamma_a\ne1\)，幂零类为 \(a\)。所有这些二面体群的导出子群都是循环群，故其二次导出子群平凡，无论是否幂零都可解。
\end{solution}

\begin{exercise}\label{b1:ex:07-central-extension}
若 \(G/Z(G)\) 幂零类至多为 \(c\)，证明 \(G\) 幂零类至多为 \(c+1\)。举例说明把 \(Z(G)\) 换成任意交换正规子群会使结论失效。
\end{exercise}
\begin{solution}
对满射 \(\pi:G\to G/Z(G)\)，式\eqref{b1:eq:lower-central-functoriality}给出
\[
\pi\Bigl(\gamma_{c+1}(G)\Bigr)
=\gamma_{c+1}\Bigl(G/Z(G)\Bigr)=1.
\]
后一等号正是商群幂零类至多为 \(c\) 的假设。因此 \(\gamma_{c+1}(G)\le\ker\pi=Z(G)\)，再与 \(G\) 取交换子，得到 \(\gamma_{c+2}(G)\le\Bigl[G,Z(G)\Bigr]=1\)。反例取 \(G=S_3\)、\(N=A_3\)：\(N\cong C_3\) 与 \(G/N\cong C_2\) 均为交换群，但 \(S_3\) 的三个二阶子群都是 Sylow \(2\) 子群且互不相同，\cref{b1:thm:finite-nilpotent-sylow}因而排除其幂零性。
\end{solution}

\begin{exercise}\label{b1:ex:07-order-p2q}
设 \(p,q\) 是不同素数。确定阶为 \(p^2q\) 的全部幂零群，直到同构。
\end{exercise}
\begin{solution}
由\cref{b1:thm:finite-nilpotent-sylow}，所求群是 Sylow 子群的直积 \(P\times Q\)，其中 \(|P|=p^2,|Q|=q\)。素数阶群循环，故 \(Q\cong C_q\)；\cref{b1:cor:order-p-square-abelian}则说明 \(P\) 交换。若 \(P\) 有 \(p^2\) 阶元，该元素生成全群，所以 \(P\cong C_{p^2}\)。否则，由\cref{b1:cor:finite-order-divides}，每个非单位元都为 \(p\) 阶。取 \(a\ne1\) 及 \(b\notin\langle a\rangle\)，两个 \(p\) 阶循环子群的交只能为 \(1\)：非平凡交会生成双方，从而使它们相等。它们交换，故\cref{b1:prop:internal-direct-product}给出乘积子群同构于 \(C_p\times C_p\)；其阶已为 \(p^2\)，因此等于 \(P\)。答案为 \(C_{p^2}\times C_q\) 与 \(C_p\times C_p\times C_q\)，两者都由有限 \(p\) 群的直积组成，按同一 Sylow 判据确实幂零。两者不同构，因为第一种有 \(p^2\) 阶元，第二种任意元素的阶都整除 \(pq\)。
\end{solution}

\begin{exercise}\label{b1:ex:07-normal-symmetric}
利用交错群单性，证明当 \(n\ge5\) 时，\(S_n\) 的正规子群只有 \(1,A_n,S_n\)。
\end{exercise}
\begin{solution}
设 \(N\trianglelefteq S_n\)，则 \(N\cap A_n\trianglelefteq A_n\)；\cref{b1:thm:alternating-simple}说明这个交为 \(1\) 或 \(A_n\)。若 \(A_n\le N\)，由\cref{b1:thm:sign}和\cref{b1:thm:first-isomorphism}，符号同态给出 \(S_n/A_n\cong C_2\)。再由\cref{b1:thm:subgroup-correspondence}，\(N/A_n\) 只能为这个二阶商群的单位子群或全群，所以 \(N=A_n\) 或 \(S_n\)。若 \(N\cap A_n=1\)，符号同态在 \(N\) 上的核为 \(N\cap A_n\)，故由\cref{b1:prop:kernel-fibers}单射，得到 \(|N|\le2\)。若阶为二，由正规性，唯一非单位元在全部 \(S_n\) 的共轭下固定，因而位于 \(Z(S_n)\)。但 \(Z(S_n)=1\)：若中心元 \(z\) 移动 \(a\) 到 \(b\ne a\)，取不同于 \(a,b\) 的 \(c\)，式\eqref{b1:eq:cycle-conjugation}给出 \(z(ac)z^{-1}=\Bigl(b,z(c)\Bigr)\ne(ac)\)，因为左侧换位移动 \(b\)，右侧不移动 \(b\)，矛盾。因此 \(N=1\)。
\end{solution}

\begin{exercise}\label{b1:ex:07-infinite-unitriangular}
将实单位上三角群依次嵌入为 \(U_1\subset U_2\subset\cdots\)，嵌入在右下角补单位元。证明并群 \(U=\bigcup_m U_m\) 的每个有限生成子群都幂零，但 \(U\) 不可解。
\end{exercise}
\begin{solution}
有限个生成元都落在某个 \(U_m\) 中。\cref{b1:prop:unitriangular-filtration}取 \(a=1,b=k\)，给出 \([U_m,1+J_k]\le1+J_{k+1}\)。以 \(\gamma_1(U_m)=U_m=1+J_1\) 为起点归纳，得到 \(\gamma_k(U_m)\le1+J_k\)，所以 \(\gamma_m(U_m)=1\)，即 \(U_m\) 幂零。式\eqref{b1:eq:lower-central-functoriality}的子群包含式又说明，所给生成元生成的子群也幂零。为证明 \(U\) 不可解，以 \(E_{ij}\) 记只有 \((i,j)\) 位置为一的矩阵。对 \(i<j<k\)，由 \(E_{ij}^2=E_{jk}^2=0\)、\(E_{ij}E_{jk}=E_{ik}\)、\(E_{jk}E_{ij}=0\)，展开四个因子得到
\[
[1+E_{ij},1+E_{jk}]=1+E_{ik}.
\]
归纳可知 \(U^{(r)}\) 含有全部 \(1+E_{i,i+2^r}\)：\(r=0\) 时这些都是 \(U\) 的元素；若结论在第 \(r\) 步成立，则 \(1+E_{i,i+2^r}\) 和 \(1+E_{i+2^r,i+2^{r+1}}\) 都在 \(U^{(r)}\) 中，其交换子由刚算出的恒等式等于 \(1+E_{i,i+2^{r+1}}\)，所以属于 \(U^{(r+1)}\)。对每个 \(r\) 取足够大的矩阵阶数，即可得到这样的非单位元，故每个导出子群非平凡，\(U\) 不可解。
\end{solution}
